%%---------- Samenvatting -----------------------------------------------------
% De samenvatting in de hoofdtaal van het document

\chapter*{\IfLanguageName{dutch}{Samenvatting}{Abstract}}

In de huidige datagedreven bedrijfswereld is een accuraat en uniform overzicht van bedrijfsdata cruciaal. IPCOM, een internationale groep actief in technische isolatie, kampt met gefragmenteerde en inconsistente data die verspreid is over diverse lokale Enterprise Resource Planning (ERP)-systemen. Het gebrek aan een centrale \textit{single source of truth} voor artikel- en leveranciersdata bemoeilijkt niet alleen de overkoepelende rapportage, maar ook de strategische besluitvorming. 

Commerciële Master Data Management (MDM)-platformen, zoals Profisee, bieden hiervoor een oplossing. Deze pakketten brengen echter aanzienlijke licentiekosten en architecturale complexiteit met zich mee, voornamelijk omdat bedrijfsdata veelal naar een extern systeem verplaatst moet worden. Dit onderzoek vertrekt daarom vanuit de volgende onderzoeksvraag: \textit{In hoeverre kan een in-house MDM-oplossing, gebouwd binnen de Microsoft Fabric-omgeving, een accuraat en schaalbaar alternatief bieden voor gespecialiseerde commerciële MDM-pakketten?}

Om deze vraag te beantwoorden, is een Proof of Concept (PoC) ontwikkeld volgens de medallion-architectuur binnen Microsoft Fabric. De dataset besloeg ruim 150.000 ruwe artikelrecords, afkomstig van IPCOM-entiteiten uit Oostenrijk, Zwitserland en Noorwegen. Na extractie is de data eerst gestandaardiseerd naar een \textit{canonical model}. Vervolgens zijn de records ontdubbeld middels een combinatie van deterministische matching (harde regels) en probabilistische fuzzy matching (op basis van de Levenshtein-afstand) via PySpark. Tot slot zijn via specifieke \textit{survivorship rules} geconsolideerde \textit{golden records} gegenereerd zonder dataverlies voor de lokale entiteiten.

De resultaten tonen aan dat de in-house ontwikkelde Fabric-oplossing zeer competitief is. Het systeem realiseerde een algehele deduplicatiegraad van 33,4\%, wat nagenoeg gelijkwaardig is aan de benchmark van 35,1\% behaald door het commerciële Profisee-platform. De deterministische regels identificeerden met succes de meerderheid van de duplicaten, terwijl de fuzzy matching cruciaal bleek om de complexere, lokaal afwijkende records te koppelen. 

Uit het onderzoek wordt geconcludeerd dat Microsoft Fabric een krachtig, naadloos integreerbaar en rendabel alternatief vormt voor externe MDM-systemen. De voornaamste beperking van de huidige PoC is het ontbreken van een grafische \textit{low-code} interface voor niet-technische \textit{data stewards}. Daarom beveelt dit onderzoek verdere integratie met tools zoals Power Apps aan als beheerlaag, evenals het onderzoeken van automatische vertaalservices om taalkundige barrières tijdens het matchen in de toekomst nog effectiever te overbruggen.